\section{Discussion}
\label{sec:discussion}

\subsection{Calibration of $\varepsilon$}

The threshold $\varepsilon_g$ in Definition~\ref{def:eps-stat} is a
policy parameter, not a protocol constant.  Its calibration involves a
tension between two goals: a small $\varepsilon$ catches even minor
statistical distortions but may block legitimate erasure requests when
the log is small; a large $\varepsilon$ permits more deletions but
weakens the guarantee against selective erasure.

We do not prescribe a universal value.  Instead, we observe that
$\varepsilon$ should be set by the regulatory authority or auditor as
part of the fairness policy $\Pi$, informed by the log size $N$, the
number of protected groups $|\mathcal{G}|$, and the baseline approval
rates.  A natural starting point is the $4/5$ths rule used in
disparate-impact analysis (EEOC, 1978): set $\varepsilon_g$ such that
a redaction batch cannot move the approval-rate ratio across the $80\%$
threshold.  Formalising this connection between $\varepsilon$ and
established legal standards is future work.

\subsection{Scalability to Large Logs}

The dominant cost of the ZK circuit is the Merkle path verification
(Section~\ref{sec:zk}), which scales as $O(B \cdot D)$ constraints
where $B$ is the batch size and $D$ is the tree depth.  For
$B = 1{,}000$ and $D = 20$ (approximately $10^6$ entries), the
estimated constraint count is $\sim$5.2M, yielding proving times of
$\sim$2--3\,s on current hardware with Barretenberg.

For logs exceeding $10^6$ entries, three mitigation strategies apply:
(1)~increasing tree depth to $D = 30$ (supporting $\sim 10^9$ entries)
adds $\sim$50\% to proving time but remains under 5\,s;
(2)~batching redactions amortises the per-redaction cost; and
(3)~recursive proof composition (e.g., IVC via Nova~\cite{nova2022})
could reduce verification to constant time regardless of batch size.
Empirical validation at scale is future work.

\subsection{Composition with Papers~1 and~2}

The three papers form a layered accountability stack:

\begin{center}
\small
\begin{tabular}{lll}
\toprule
\textbf{Paper} & \textbf{Layer} & \textbf{Predicate} \\
\midrule
1 (BTV)   & Construction  & $V \multimap (E \otimes C)$ \\
2 (BTV-T) & Persistence   & $\mathsf{DeliveryToken} \multimap (V \otimes R)$ \\
3 (This)  & Redaction     & $\mathsf{Valid}(L_t \xrightarrow{R} L_{t+1})$ \\
\bottomrule
\end{tabular}
\end{center}

\noindent Papers~1 and~2 enforce invariants at compile time via linear
resource types in Rust.  Paper~3 operates at the cryptographic layer
and is checked at runtime by a SNARK verifier.  The layers are
complementary: the type system ensures no decision escapes without a
persistence record; the SNARK verifier ensures no persistence record
can be selectively erased without statistical justification.  Neither
mechanism subsumes the other.

A system implementing all three layers guarantees: (a)~no decision
without evidence, (b)~no delivery without verifiable persistence, and
(c)~no redaction without statistical consistency.  Formal composition
of the three predicates into a single end-to-end security theorem is
future work.

\subsection{Privacy Leakage from Commitment Transitions}

As noted in Section~\ref{sec:protocol}, each redaction publishes a
transition from $\mathsf{Com}(S(L_t); \rho_t)$ to
$\mathsf{Com}(S(L_{t+1}); \rho_{t+1})$.  An observer who sees $k$
successive transitions learns $k$ equations in the group counters
$(s_g, n_g)$.  For small groups, this constitutes $O(\log N)$ bits of
leakage per redaction.

Two mitigations bound this leakage: (1)~batching multiple redactions
into a single transition reduces the number of observable transitions;
(2)~adding differential-privacy noise to the committed statistics
(at the cost of weakening the $\varepsilon$ guarantee by a controlled
amount).  Quantifying the privacy--accuracy trade-off under
differential privacy composition is future work.

\subsection{Limitations}

\begin{itemize}
  \item The protocol assumes a trusted setup for the SNARK
  (Groth16/PLONK).  Universal setups (e.g., KZG ceremony) mitigate
  but do not eliminate this assumption.
  \item The fairness policy $\Pi$ must be agreed upon before the log
  begins; changing $\Pi$ mid-stream requires a migration protocol not
  defined here.
  \item The circuit operates on committed demographic attributes.  If
  the operator commits false attributes at ingestion time, the
  statistical guarantees are vacuous.  This is analogous to the
  semantic-truthfulness limitation of Paper~1
  (Section~6.1 of~\cite{btv2026}).
  \item Proving time scales linearly with batch size.  For real-time
  redaction of individual entries, the $\sim$2\,s overhead may be
  acceptable; for bulk GDPR compliance sweeps affecting thousands of
  entries, parallelisation or recursive composition is needed.
\end{itemize}
